% ID: SEL-D0002-OA-E03
% ESTADO: validada. S003, 2026-11-07; tres pasadas del mismo autor, comprobación separada y revisión didáctica/visual.
\textbf{Datos.}\par\nopagebreak[4]
Sea $V$ el suceso ``presenta avería''. Los modelos $A,B,C$ forman una partición: $P(A)=0.25$, $P(B)=0.40$ y $P(C)=1-0.25-0.40=0.35$. Las tasas de avería son probabilidades condicionadas, no probabilidades conjuntas.

\textbf{a) Probabilidad total.}\par\nopagebreak[4]
\[
\begin{aligned}
P(V)&=P(A)P(V\mid A)+P(B)P(V\mid B)+P(C)P(V\mid C)\\
&=0.25(0.15)+0.40(0.10)+0.35(0.12)\\
&=0.0375+0.04+0.042=\boxed{0.1195}.
\end{aligned}
\]
Esto equivale al $11.95$\,\%.

\textbf{b) Complementario dentro del modelo A.}\par\nopagebreak[4]
Ya sabemos el modelo: no multiplicamos de nuevo por $P(A)$.
\[
P(V^c\mid A)=1-P(V\mid A)=1-0.15=\boxed{0.85}.
\]

\textbf{c) Unión.}\par\nopagebreak[4]
Los sucesos $V$ y $C$ no son incompatibles. Aplicamos inclusión-exclusión:
\[
P(V\cup C)=P(V)+P(C)-P(V\cap C)
=0.1195+0.35-0.042=\boxed{0.4275}.
\]
Comprobación por sucesos disjuntos: todos los patinetes C, más los averiados A y B, suman $0.35+0.0375+0.04=0.4275$. Además, las probabilidades conjuntas de avería y no avería suman 1.
